% ===== 问题二正文 =====
% 问题二正文；三张示意图的 TikZ 宏在 figures/q2_schematics.tex 中定义。
\subsection{问题二：第二检测点选择模型的建立与求解}\label{sec:q2}

问题二中只进行了一次测向，需要确定第二个检测点的位置。此时干扰源的位置和接收半径均未知，第二检测点的优劣只能依据首测信息判断。本节先在模型准备中给出首测后源的可能范围，并分析第二次测向的几何作用；再以第二检测点为决策变量建立优化模型；然后给出模型的求解方法；最后用实际观测数据检验。本节建立的单源定位模型在第三问中反复调用。

\subsubsection{模型准备}\label{sec:q2-1}

\textbf{（1）符号约定。}记所考察的频道为 $f$，源的位置为 $p$，首个检测点为 $x_1$，测得的示向度为 $\theta_1$，示向方向的单位向量记为 $u(\theta_1)=(\cos\theta_1,\sin\theta_1)$。场地为半径 $1800$ m 的圆，记作 $\Omega$；$B(x,r)$ 表示以 $x$ 为圆心、半径为 $r$ 的圆盘。接收半径 $R_f$ 在一次任务中固定，仅知其取值范围为 $1000$ 至 $1500$ m。测角误差按题设取 $\pm\varepsilon$，$\varepsilon=1^\circ$。

\textbf{（2）首测后的定位区域。}首测得到示向度 $\theta_1$ 后，源的位置 $p$ 需同时满足三个条件：

\begin{itemize}
\item \textbf{角度条件：}$p$ 位于示向度 $\theta_1\pm1^\circ$ 所夹的楔形 $W(x_1,\theta_1)$ 内，不等式写法同第一问的式\eqref{eq:q1-1}；
\item \textbf{距离条件：}在 $x_1$ 收到了信号，故 $\|p-x_1\|\le1500$ m；
\item \textbf{场地条件：}$p\in\Omega$。
\end{itemize}

三者的交集即首测后源可能所在的范围，沿用第一问的叫法称为定位区域：

\begin{equation}\label{eq:q2-1}
C_f^{(1)}=\Omega\cap B(x_1,1500)\cap W(x_1,\theta_1).
\end{equation}

\textbf{（3）定位区域的多边形近似。}$C_f^{(1)}$ 的边界含有两段圆弧，不便像第一问那样用顶点计算。将两个圆分别换成其外切正 $128$ 边形，得到凸多边形 $\widehat C_f^{(1)}\supseteq C_f^{(1)}$。采用外切而非内接多边形，是为了使近似只扩大、不缩小区域，真实源不会因近似而被排除。后文的选点和清除判断均以 $\widehat C_f^{(1)}$ 为准。另记 $\rho(C)$、$z(C)$ 分别为区域 $C$ 的最小外接圆半径和圆心，第三问中将该圆简称为包围圆。

\textbf{（4）第二次测向的几何作用。}首测后的定位区域是一条张角 $2^\circ$、长达 $1500$ m 的窄带：角度已经足够准，缺的是距离。第二次测向要做的，是用第二条视线把这条窄带截断，截断的效果取决于两条视线的夹角，即测量学中的交会角。图\ref{fig:q2-cross}比较了两种方案：

\begin{itemize}
\item \textbf{沿示向方向前进再测：}两条视线几乎平行，只截去窄带靠近 $x_1$ 的一段，剩余部分仍然狭长；
\item \textbf{横向偏移后再测：}两条视线斜交，交集只是一个小四边形。
\end{itemize}

\begin{center}
\begin{minipage}{\linewidth}
\centering
\QTwoCrossFigure
\captionof{figure}{两种第二检测点的比较。沿示向方向前进只能截短楔形；横向偏移后两条视线斜交，定位区域才明显缩小。}
\label{fig:q2-cross}
\end{minipage}
\end{center}

两种方案的差别可以用交会角定量刻画。设两次检测点到源的距离分别为 $d_1$、$d_2$，交会角为 $\gamma$。每个楔形在源附近的宽度约为 $2d_i\varepsilon$（$\varepsilon$ 按弧度计），两条窄带斜交所得四边形的对角线长度约为

\begin{equation}\label{eq:q2-2}
D\approx\frac{2(d_1+d_2)\,\varepsilon}{\sin\gamma}.
\end{equation}

由式\eqref{eq:q2-2}，$\gamma$ 越接近 $90^\circ$，定位区域越小；$d_2$ 越小，定位区域也越小。因此第二检测点应当靠近源，并与首测视线拉开一定的横向距离。

\textbf{（5）接收范围的限制。}横向偏移不能没有限度。示向度只有在源的接收半径之内才能测得，而接收半径未知，只知道不小于 $1000$ m。第二检测点若离源过远，将收不到信号，移动和检测的时间随之浪费。因此第二检测点的选择需要兼顾两方面：增大交会角以缩小定位区域，保证接收以免测向失败。前者构成下面模型的目标函数，后者构成约束条件。

\subsubsection{模型的建立}\label{sec:q2-2}

\textbf{（1）决策变量。}待确定的是第二检测点的位置，记为 $x\in\mathbb R^2$。为使模型在第三问中同样适用，记机器狗当前所在位置为 $x_0$：第二问中 $x_0=x_1$，第三问中为当时的实时位置；当前的定位区域记为 $\widehat C_f$，第二问中即 $\widehat C_f^{(1)}$。

\textbf{（2）约束条件。}第一条约束是保证接收。第二次测向若收不到信号，便无法形成双站交会。接收半径虽然未知，但不小于 $1000$ m；因此只要 $x$ 到源的每一个可能位置的距离都不超过 $1000$ m，则无论源的实际位置和接收半径如何，到达 $x$ 后必定能收到信号。满足该条件的位置集合称为可靠接收区，记作 $Q_f$。该条件涉及区域内的无穷多个点，但由于区域是凸多边形，只需检验其顶点：

\begin{equation}\label{eq:q2-3}
Q_f=\bigl\{x:\ \|x-a\|\le1000\ \text{对}\ \widehat C_f\ \text{的每个顶点}\ a\ \text{成立}\bigr\}.
\end{equation}

理由与第一问求直径时相同：多边形内任一点都是顶点的凸组合，它到 $x$ 的距离不超过各顶点到 $x$ 距离的最大值；而顶点本身属于区域。因此到所有顶点的距离不超过 $1000$ m，等价于到区域内所有点的距离不超过 $1000$ m，检验一个候选点只需计算与顶点个数相同的若干次距离。$Q_f$ 是有限个圆盘的交集，因而是凸集。

式\eqref{eq:q2-3}便于计算，但没有直接给出该区域的位置和大小，下面给出它的一个可以直接写出的子区域。记 $\widetilde R$ 为 $\widehat C_f^{(1)}$ 的顶点到 $x_1$ 的最大距离（略大于 $1500$ m），则整个定位区域落在以 $x_1$ 为顶点、半径 $\widetilde R$、半角 $\varepsilon$ 的扇形内。过 $x_1$ 和扇形的两个远端点作圆，其圆心 $z_0$ 位于示向线上，半径为

\begin{equation}\label{eq:q2-4}
r_0=\frac{\widetilde R}{2\cos\varepsilon}\approx750\ \text{m},\qquad z_0=x_1+r_0\,u(\theta_1).
\end{equation}

该圆恰好覆盖整个扇形：扇形内任一点到 $z_0$ 的距离平方是它到 $x_1$ 距离的凸函数，最大值只能在扇形顶点或远端弧上取得，而这两处到 $z_0$ 的距离均不超过 $r_0$。于是，以 $z_0$ 为圆心、半径 $1000-r_0\approx250$ m 的圆盘内任一点 $x$，到任一可能源 $p$ 的距离满足 $\|x-p\|\le\|x-z_0\|+\|z_0-p\|\le1000$，即该圆盘整体包含于 $Q_f$。换言之，沿示向方向前进约 $750$ m，以该处为圆心、半径约 $250$ m 的圆盘内的任一位置，都能保证再次收到信号。场地边界只会缩小定位区域、扩大 $Q_f$，因此上述结论对任意首测位置均成立。图\ref{fig:q23-1}按实际比例画出了首测楔形、覆盖圆和可靠接收区。

\begin{center}
\begin{minipage}{\linewidth}
\centering
\QTwoRegionFigure
\captionof{figure}{首测楔形、覆盖圆与可靠接收区（按实际比例绘制）。蓝色为可靠接收区 $Q_f$ 的一部分，由包含楔形的三角形的三个顶点按式\eqref{eq:q2-3}画出；虚线圆为覆盖整个楔形的圆 $B(z_0,r_0)$；绿色为以 $z_0$ 为圆心、半径约 $250$ m 的圆盘，整体落在 $Q_f$ 内。}
\label{fig:q23-1}
\end{minipage}
\end{center}

式\eqref{eq:q2-3}只是保证接收的充分条件，并非能够接收的最大范围。首测收到信号还蕴含 $R_f\ge\|p-x_1\|$，若将这一信息一并利用，可以进一步放宽部分位置。本文统一采用 $1000$ m 这一界限，一方面计算简单，另一方面在第三问的多次测向中可以直接沿用。

第二条约束是不重复检测：$x$ 不能取该频道已经测过的位置。题设指出同一地点重复检测不能消除误差，重复检测只会增加耗时。

\textbf{（3）目标函数。}选点的目的是减少任务耗时，因此以前往某点测向一次所引起的各项时间之和作为代价。测向后定位区域收缩的程度取决于源的真实位置，而后者未知。本文的处理办法是假设源位于定位区域内的若干代表位置上：区域中心 $z=z(\widehat C_f)$，以及从顶点中均匀选取的至多 $4$ 个点，各自向中心移动四分之一（选取靠近顶点的位置是为了兼顾源处于区域边缘的情形，向内移动则避免代表位置恰好落在边界上），记为 $p^{(1)},\dots,p^{(h)}$，$h\le5$。对候选点 $x$ 和假想的源位置 $p^{(k)}$，由 $x$ 指向 $p^{(k)}$ 的方向即为一条不含误差的示向度，将其作为新的测向结果，用第一问的方法在区域上增加一个楔形约束，求出新区域的包围圆半径，记为 $r_k(x)$。若 $x$ 与 $p^{(k)}$ 的距离不超过 $5$ m，按题设设备将给出近场反馈，可直接在该处清除。

移动速度 $v=5$ m/s，可靠清除的半径门限取 $r_c=19.9$ m（略小于 $20$ m 的光学清除距离）。候选点 $x$ 的代价定义为

\begin{equation}\label{eq:q2-5}
J(x)=\frac{\|x-x_0\|}{v}
+\frac1h\sum_{k=1}^{h}\left[\frac{\|x-p^{(k)}\|}{v}+P_k(x)\right].
\end{equation}

三项依次为：由当前位置移动到候选点的时间；测向后由候选点移动到源的时间，对 $h$ 个假想位置取平均；测向后仍未达到清除条件的罚时 $P_k(x)$。罚时的取法为：若 $r_k(x)$ 不超过 $r_c$，或已得到近场反馈，则罚时为零；否则至少还需一次检测和一次换频，计 $6$ s，半径每超出门限 $1$ m 再加 $0.1$ s，以反映后续额外的移动。换频与本次检测的时间对所有候选点相同，不影响排序，故不计入。

\textbf{（4）模型汇总。}综上，第二检测点的选择模型为

\begin{equation}\label{eq:q2-6}
\begin{aligned}
\min_{x}\quad & J(x)\\
\text{s.t.}\quad & \|x-a\|\le1000,\quad a\ \text{为}\ \widehat C_f\ \text{的每个顶点},\\
& x\ \text{不是该频道已测过的位置}.
\end{aligned}
\end{equation}

约束条件保证第二次测向必定收到信号，目标函数在移动路程与测后剩余区域之间取得平衡。等权的假想位置仅用于比较候选点，并不代表源位置的概率分布，因此模型\eqref{eq:q2-6}给出的是一条可执行的选点规则，而非所有情形下的最优解。实际到达第二检测点并完成测向后，只用真实示向度更新定位区域，假想位置不再参与。

\subsubsection{模型的求解}\label{sec:q2-3}

模型\eqref{eq:q2-6}的可行域 $Q_f$ 是连续区域，目标函数中又含有区域求交和最小外接圆的计算，难以求得解析解。考虑到实际只需要一个足够好的检测点，本文采用有限候选点枚举的方法求解：在可行域内布设若干具有代表性的候选点，逐一计算代价，取代价最小者。

\textbf{（1）候选点的生成。}候选点分两组，布设方式见图\ref{fig:q2-cand}。

第一组沿前进方向布设。在 $x_0$ 到区域中心 $z$ 的连线上取中点和终点，各自向两侧横向偏移 $50$ m 和 $150$ m，连同中心 $z$ 共 $9$ 个点。这一组保证候选中总有靠近区域、路程较短的位置。

第二组沿定位区域的长轴布设。由式\eqref{eq:q2-2}，最需要截断的是区域最长的方向。取多边形中相距最远的一对顶点，其连线方向为长轴；将全部顶点投影到长轴上，得到区域沿长轴的长度 $w$。在长轴的 $1/4$、$3/8$、$5/8$、$3/4$ 四个分位点上，各取轴上一点及两侧横向偏移 $b_*$ 的两点，再在中心 $z$ 两侧各取一个横向偏移 $b_*$ 的点，共 $14$ 个点。横向偏移量随区域大小调整：

\begin{equation}\label{eq:q2-7}
b_*=\min\{200,\ \max\{20,\ w/8\}\}\ \text{m},
\end{equation}

即取区域长度的八分之一，并限制在 $20$ m 至 $200$ m 之间。区域越长横向偏移越大，上限 $200$ m 用于避免候选点超出可靠接收区。分位点中加入靠内的 $3/8$ 和 $5/8$，是因为首测楔形很长时，$1/4$ 和 $3/4$ 处的横向偏移点往往已超出可靠接收区，靠内的分位点可以保留横向选择。

两组候选点合并后，去除重复点、该频道已测过的位置以及不满足式\eqref{eq:q2-3}的点，至多剩余 $23$ 个。

\begin{center}
\begin{minipage}{\linewidth}
\centering
\QTwoCandidateFigure
\captionof{figure}{两组候选检测点的布设示意。第一组沿当前位置到区域中心的连线布点，第二组沿定位区域的长轴布点；超出可靠接收区的点剔除。}
\label{fig:q2-cand}
\end{minipage}
\end{center}

\textbf{（2）代价的计算。}对每个候选点按式\eqref{eq:q2-5}计算代价，取代价最小者作为第二检测点。第三问中同一个源可能需要多次测向，若定位区域收缩后已无可行的新候选点，则改用区域中心及其附近的偏移点，再无新点时转入第三问的光学网格清除。

\textbf{（3）求解步骤。}求解流程见图\ref{fig:q2-flow}，步骤归纳为算法\ref{alg:q2-probe}。

\begin{center}
\begin{minipage}{\linewidth}
\centering
\QTwoFlowFigure
\captionof{figure}{第二检测点选择模型的求解流程。第二问执行到第一次更新定位区域为止；第三问中同一个源按此流程反复测向，直至达到清除门限。}
\label{fig:q2-flow}
\end{minipage}
\end{center}

\par\medskip
\noindent\begin{minipage}{\linewidth}
\small
\refstepcounter{algorithm}\label{alg:q2-probe}
\textbf{算法\thealgorithm：第二检测点选择模型的求解}\par\smallskip
\noindent\textbf{输入：}当前位置 $x_0$、首测示向度 $\theta_1$、定位区域 $\widehat C_f$、该频道已测过的位置
\begin{enumerate}
\setlength{\itemsep}{2pt}
\setlength{\parskip}{0pt}
\item 求 $\widehat C_f$ 的包围圆圆心 $z$ 和长轴，生成两组候选点；
\item 去重，去除已测过的位置，按式\eqref{eq:q2-3}剔除不在可靠接收区内的点；
\item 选取假想源位置，对每个候选点按式\eqref{eq:q2-5}计算代价 $J(x)$；
\item 取代价最小的点作为第二检测点；
\item 前往该点实际测向，用真实示向度更新 $\widehat C_f$；得到近场反馈则直接清除。
\end{enumerate}
\end{minipage}
\par\medskip

\subsubsection{结果分析}\label{sec:q2-4}

以第三问随机场景 800035 中频道 $20$ 的实际运行数据检验模型。该频道首测在原点，第二检测点由算法\ref{alg:q2-probe}选出，两次测向的数据及定位区域的变化见表\ref{tab:q2-example}。

\Needspace{38mm}
\begin{center}
\begin{minipage}{\linewidth}
\centering\small
\captionof{table}{频道 20 的两次测向与定位区域的变化}
\label{tab:q2-example}
\vspace{4pt}
\begin{tabular}{@{}lcrr@{}}
\toprule
观测 & 检测点坐标 / m & 示向度 & 定位区域包围圆半径 / m \\
\midrule
原点首测 & $(0.000,\ 0.000)$ & $125.00^\circ$ & $750.226$ \\
第二次测向 & $(-170.646,\ 568.027)$ & $148.14^\circ$ & $65.012$ \\
\bottomrule
\end{tabular}
\end{minipage}
\end{center}

算例结果可归纳为三点。

\textbf{第二检测点满足接收约束。}模型选出的第二检测点沿示向方向前进约 $563$ m、横向偏移约 $186$ m，落在可靠接收区内：它到首测定位区域各顶点的最大距离为 $960$ m，小于 $1000$ m。实际测向也再次收到了示向度。

\textbf{定位区域显著缩小。}两条示向线的交会角约 $23^\circ$，一次测向即将定位区域的包围圆半径由 $750$ m 缩小到 $65$ m，缩小幅度超过 $90\%$。作为对比，若沿示向方向直行同样的距离再测，两条示向线近乎平行，定位区域仍将长达数百米。

\textbf{结果与理论估计一致。}由两条示向线的交点估计，源到两个检测点的距离分别约为 $1000$ m 和 $470$ m，代入式\eqref{eq:q2-2}得 $D\approx130$ m，与包围圆直径 $2\times65=130$ m 相符，说明模型准备中的交会误差估计可以用于预判选点效果。

此时 $65$ m 尚未达到 $19.9$ m 的清除门限，在第三问中该源需按同一模型再测向一次；选点规则对任务总耗时的影响在第三问的对比实验中评价。

\subsubsection{问题二的结论}\label{sec:q2-5}

综合以上分析，问题二的回答如下：

\begin{itemize}
\item \textbf{候选区域。}第二检测点应选在可靠接收区 $Q_f$ 内，即到首测定位区域各顶点的距离均不超过 $1000$ m 的位置。无论首测位置如何，该区域都包含一个可直接确定的圆盘：沿示向方向前进约 $750$ m，以该处为圆心、半径约 $250$ m。
\item \textbf{选择策略。}在可靠接收区内不应沿示向线直行，而应横向偏移，使两条示向线形成尽可能大的交会角，区域越长横向偏移越大。具体位置由模型\eqref{eq:q2-6}在两组候选点中比较移动路程与测后剩余区域，取代价最小者。
\end{itemize}

该策略保证真实源始终位于定位区域内、第二检测点必定收到信号、清除判定可靠；其对任务耗时的实际改善在第三问的完整任务中检验。
